\documentclass[../DoAn.tex]{subfiles}
\begin{document}

Chương này trình bày một số lưu ý quan trọng khi sử dụng tài liệu tham khảo trong báo cáo đồ án tốt nghiệp.

Về nguyên tắc trích dẫn, sinh viên cần tuân thủ nghiêm ngặt quy tắc không được sao chép nguyên văn tài liệu của người khác mà không ghi rõ nguồn. Mọi hình vẽ, bảng biểu, công thức và dữ liệu không phải do sinh viên tự tạo ra đều phải có trích dẫn nguồn gốc rõ ràng. Trong văn phong khoa học, sinh viên tuyệt đối không được sử dụng các nguồn như Wikipedia hoặc bài giảng chưa được công bố chính thức (lecture slides) làm tài liệu tham khảo.

Về định dạng trích dẫn, đồ án sử dụng định dạng trích dẫn theo chuẩn IEEE được tích hợp trong template. Các tài liệu tham khảo được liệt kê trong danh sách theo thứ tự xuất hiện trong văn bản, sử dụng hệ thống đánh số tự động. Khi trích dẫn, sinh viên sử dụng lệnh \verb|\cite{key}| trong LaTeX, trong đó key là khóa được định nghĩa trong file \verb|Danh_sach_tai_lieu_tham_khao.bib|.

Về các loại tài liệu tham khảo được chấp nhận, các bài báo khoa học (journal articles) được đăng trên các tạp chí uy tín cần được ưu tiên sử dụng. Các sách chuyên ngành (books) từ các nhà xuất bản uy tín là nguồn tham khảo đáng tin cậy. Các bài báo hội nghị (conference papers) từ các hội nghị quốc tế có bình duyệt cũng được chấp nhận. Các tài liệu chính thức từ các tổ chức uy tín như tài liệu API, tài liệu chuẩn công nghệ là nguồn tham khảo tốt cho các phần công nghệ. Luận văn, luận án (theses, dissertations) từ các trường đại học uy tín có thể được sử dụng khi liên quan trực tiếp đến đề tài.

\end{document}
